
\section{Related Work}

\subsection{Low-Rank Matrix Factorization}

\subsubsection{Classical SVD}

Eckart \& Young (1936) proved that the optimal rank-$r$ Frobenius-norm 
approximation of $W$ is the truncated SVD $W_r = U_{:, :r} \Sigma_{:r} V_{:r, :}^\top$. 
This is the foundation for all SVD-based compression.

\subsubsection{Randomized SVD}

Halko et al. (2011) provide a randomized algorithm for truncated SVD with 
complexity $\mathcal{O}(mn \log r)$, trading accuracy for speed. For our 
layer sizes (up to 3072$\times$768), exact SVD is fast enough.

\subsubsection{Weighted Low-Rank Approximation}

Golub et al. (1987), Srebro \& Jaakkola (2003). Weighted approximation 
minimizes $\|(W - W_r) D\|_F$ where $D$ is a weighting matrix. ASVD uses 
$D = \text{diag}(\sqrt{F})$ where $F$ is the activation energy.

\subsubsection{Tensor Decomposition}

Tensor-train (Novikov et al., 2015) uses higher-order factorization. 
ASVD is a matrix method.

\subsection{Activation-Aware Compression}

\subsubsection{ASVD Paper}

Yuan et al. (2023) -- arXiv:2307.07382 -- introduced ASVD for LLM 
compression. The core idea: scale the weight matrix by the input activation 
magnitude before SVD. Our implementation extends this with explicit $\Sigma$, 
energy-based dynamic rank, and pre-flight validation.

\subsubsection{AWQ}

Lin et al. (2023) use activation magnitude for weight quantization. 
ASVD uses it for low-rank factorization.

\subsubsection{SmoothQuant}

Xiao et al. (2023) migrate activation scale to weights for INT8 quantization. 
ASVD performs an analogous migration but for the SVD input space.

\subsubsection{A3}

Our companion module A3 uses diagonal Fisher weighting with LayerNorm 
recalibration. ASVD differs by keeping $\Sigma$ explicit and using energy-based 
rank selection instead of adaptive retry.

\subsubsection{AASVD}

Our companion module AASVD uses full covariance weighting with streaming 
EMA. ASVD uses a simpler diagonal weighting.

\subsection{Parameter-Efficient Fine-Tuning}

\subsubsection{LoRA}

Hu et al. (2021) -- low-rank adapters with $B=0$ initialization. ASVD uses 
LoRA for recovery after compression.

\subsubsection{AdaLoRA / DoRA}

Zhang et al. (2023), Liu et al. (2024) -- adaptive rank allocation and 
weight decomposition. ASVD uses standard LoRA.

\subsection{Position of Our Work}

\begin{table}[H]
\centering
\small
\begin{tabular}{lll}
\toprule
\textbf{Axis} & \textbf{Prior work} & \textbf{ASVD} \\
\midrule
\textbf{Weighting} & Unweighted or diagonal & Diagonal (activation energy) \\
\textbf{$\Sigma$ handling} & Implicit or folded & \textbf{Explicit matrix} \\
\textbf{Rank selection} & Fixed or adaptive-retry & \textbf{Energy-based dynamic} \\
\textbf{Validation} & Rare & \textbf{Strict Frobenius with abort} \\
\textbf{Pre-flight check} & None & \textbf{Extraction validation} \\
\textbf{Multi-architecture} & Usually one family & Four families \\
\bottomrule
\end{tabular}
\end{table}

\textbf{Key distinction.} ASVD's explicit $\Sigma$ matrix distinguishes it from 
A3 and AASVD, which fold the singular values into the factors. This enables 
inspection and rank reduction without re-factorization.

\subsection{Comparison to Companion Modules}

\begin{lstlisting}

---
\end{lstlisting}